\section{不同空间上的映射与比较}
\label{sec:measure-cross-space}

本节研究空间之间的质量怎样传递、联合记录和比较。
确定映射为每个输入指定一个目标点；\term{测度核}为每个输入指定目标空间上的一项非负测度。
二者在相应的可测和质量条件下产生输出测度，再按质量、\term{绝对连续}性或矩条件核对上一节各项比较的适用范围。
联合测度保留输入与输出之间的配对关系：在有限输入测度和核有统一质量界的范围内可以正向构造联合测度；从联合测度反求核则要另查空间条件。
若直接比较原空间之间的配对代价，还要另外给定跨空间成本。
各小节按这些操作各自回答的问题分工，再按定义、运算和比较的依赖安排顺序；输出匹配程度与实现对应所需的代价分别判断。

\subsection{\term{可测映射}的推前与密度变换}

确定映射先回答目标集合接收多少输入质量，再由输出积分公式判断密度怎样变化。
推前不以映射可逆或具有密度为前提；一般定义和复合在本小节集中建立，欧氏可逆换元只回用前面的微积分结论，不重新证明。

\subsubsection{\term{推前测度}与输出积分}

\paragraph{推前的定义与质量边界}

固定可测映射$T:(X,\mathcal A)\to(Y,\mathcal B)$和$X$上的测度$\mu$。
对目标可测集$B\in\mathcal B$定义
\[
 (T_\#\mu)(B)=\mu(T^{-1}(B)).
\]
可测性保证每个目标事件的原像都是$\mathcal A$中的合法集合，从而有权重；这里的$T^{-1}(B)$是集合原像，不要求$T$存在逆函数。

$T_\#\mu$确实是$Y$上的测度。
空集原像为空，质量为零；若$B_n$两两不交，原像$T^{-1}(B_n)$也两两不交且
\[
 T^{-1}\!\left(\bigcup_nB_n\right)=\bigcup_nT^{-1}(B_n),
\]
由$\mu$的可数可加性得到$T_\#\mu$可数可加。
又因为$T^{-1}(Y)=X$，总质量守恒：
\[
 (T_\#\mu)(Y)=\mu(X).
\]
所以\term{有限测度}推前后仍有限，概率推前后仍是概率。

$\sigma$有限性不能无条件保持。
实轴上的勒贝格测度经常值映射$T(x)=0$推前后得到$\delta_0$：目标点$0$的质量为$+\infty$，任何包含它的集合都不能具有有限测度，输出无法由有限质量集覆盖。
若两个带质量源点映到同一目标点，目标点收到的质量是两项权重之和；输出测度由这些集合权重完全确定，它是否在某个参考测度下有密度则要另行判断。

\paragraph{输出观测在输入端的表示}

要计算一般输出观测的积分，先从指标函数做起。
对$g=\mathbf 1_B$，
\[
 \int_Y g\,d(T_\#\mu)
  =(T_\#\mu)(B)=\mu(T^{-1}B)
  =\int_X\mathbf 1_B(Tx)\,d\mu(x)
  =\int_X g(Tx)\,d\mu(x).
\]
沿“指标—非负简单函数—非负可测函数单调逼近”的路线推广，对非负可测$g:Y\to[0,\infty]$有
\begin{equation}
 \int_Y g\,d(T_\#\mu)=\int_X g(Tx)\,d\mu(x),
 \label{eq:measure-pushforward-integral}
\end{equation}
允许两边共同为$+\infty$。
有符号$g$分别处理正负部分，只在绝对可积或能避免无穷相减的合法范围内合并。
无须为一般空间重新证明积分构造。

再取可测映射$S:(Y,\mathcal B)\to(Z,\mathcal C)$。由
\[
 (S\circ T)^{-1}(C)=T^{-1}(S^{-1}(C))
\]
直接得到
\[
 (S\circ T)_\#\mu=S_\#(T_\#\mu):
\]
先把质量送到$Y$、再送到$Z$，与直接送到最终点一致。
若某个可测映射$T:X\to X$满足$T_\#\mu=\mu$，称它\term{保持测度}，这只是输出权重与输入权重相同的特例；本章不展开遍历性或Haar测度存在理论。

\subsubsection{输出密度与质量重组}

\paragraph{回用欧氏换元计算输出密度}

式~\eqref{eq:measure-pushforward-integral}把目标观测复合到输入端，不要求输出具有密度，也不是在原空间乘一个Radon–Nikodym系数。
给定参考体积和映射的正则性后，才谈得上输出密度。

设$U,V\subseteq\mathbb R^d$是同维欧氏开集，$T:U\to V$是$C^1$微分同胚，输入测度$\mu$相对于$U$上的勒贝格体积有非负密度$p$。
这里要求全局一一对应、逆映射光滑；局部可逆不足以替代这些假设。
对非负测试函数$g$，先用推前积分恒等式，再回用多元换元公式：
\[
 \int_V g(y)\,d(T_\#\mu)(y)
  =\int_U g(Tx)p(x)\,dx
  =\int_V g(y)\,p(T^{-1}y)\,|\det DT^{-1}(y)|\,dy.
\]
由于这对所有非负$g$成立，输出密度的几乎处处表示为
\[
 q(y)=p(T^{-1}y)\,|\det DT^{-1}(y)|.
\]
利用逆导数关系$DT^{-1}(Tx)=(DT(x))^{-1}$，也可写成
\[
 q(Tx)\,|\det DT(x)|=p(x).
\]
体积伸缩因子与密度系数相乘正好保持质量；有符号测试还需积分合法性。
这里只衔接一般推前和已有换元结论，不重做微积分中的证明。
平移映射给出$q(y)=p(y-b)$；可逆线性映射$T(x)=Ax$给出
\[
 q(y)=\frac{p(A^{-1}y)}{|\det A|};
\]
负比例缩放由行列式绝对值处理体积。
多原像、降维和退化映射不满足全局微分同胚条件，必须改用原像计算，不能挑选某一个逆分支硬套公式。

\paragraph{合并、投影与原子输出}

离散合并时，目标点收到全部原像源点的权重之和；连续映射若有多个逆分支，也要把所有原像分支的贡献相加。
推前本身始终有定义，但目标密度不能由单个全局逆自动保证。

设$\mathbb R^2$上的非负可积密度为$p(x,y)$（对乘积可测结构可测；若只在勒贝格完备化下可测，先取几乎处处相同的乘积可测代表），考虑投影$T(x,y)=x$。
对一维Borel集$B$，原像为$B\times\mathbb R$，由Tonelli定理，
\[
 (T_\#\mu)(B)
  =\int_{B\times\mathbb R}p(x,y)\,dx\,dy
  =\int_B\left(\int_{\mathbb R}p(x,y)\,dy\right)dx.
\]
因此边缘密度
\[
 q(x)=\int_{\mathbb R}p(x,y)\,dy
\]
正是该投影的输出密度，目标参考是一维勒贝格测度。
这是已有累次积分的用途，不再铺开构造。

改为常值映射$T(x,y)=y_0$时，有限输入总质量$m$被推为原子$m\delta_{y_0}$，因为目标集合的原像只能是全空间或空集。
在正维欧氏目标空间中且$m>0$时，它相对于勒贝格测度没有普通密度；$m=0$则给出零测度，要与非零原子分开。
投影和常值输出说明：不可逆映射可以保持密度，也可以产生原子，两种情形都可能合并原本不同的点。
确定映射仍要求每个输入只对应一个输出；若同一输入还要把质量分给多个目标，就要由下一小节的测度值规则取代单个输出点。

\subsection{测度核与\term{测度值映射}}

本小节把“每个输入对应一个目标点”扩展为“每个输入对应目标集合上的一项非负测度”。
新对象仍是映射，但输出从点换成测度；随输入变化的集合质量必须可测，才能对输入测度累计。
先给一般测度核的定义，再为积分、复合和联合构造明确采用有限输入测度及统一有限质量界的核；这是本章选择的充分工作范围，不作为这些运算的最弱条件。
总质量为$1$是额外的归一化条件，不写入一般核的定义。

\subsubsection{随输入变化的输出权重}

\paragraph{测度核的两项条件}

先看有限输入、输出集合上的非负矩阵：固定一行，矩阵元素给出该输入对应各目标点的质量，对一组目标点求和便得到集合质量。
不要求各行之和相同或等于$1$；当输入扩展到一般集合时，还必须让每个目标集合的质量随输入可测，才能对输入测度累计。

\begin{definition}[测度核]
\label{def:measure-kernel}
设$(X,\mathcal A)$与$(Y,\mathcal B)$为可测空间。映射
\[
 K:X\times\mathcal B\to[0,\infty]
\]
若满足：
\begin{enumerate}
 \item 对每个$x\in X$，$K_x=K(x,\cdot):\mathcal B\to[0,\infty]$是$(Y,\mathcal B)$上的非负测度；
 \item 对每个$B\in\mathcal B$，函数$x\mapsto K(x,B)$关于$\mathcal A$可测，
\end{enumerate}
就称$K$为从$X$到$Y$的测度核，记作$K:X\leadsto Y$。
\end{definition}

第一项约束目标端的集合赋权，第二项约束输入端的可测依赖，缺一不可。
这里的测度核与函数空间章中表示相似度的正定核是完全不同的对象。

一般核定义本身没有质量上界。
本章主要运算进一步要求存在与$x$无关的常数$M<\infty$，使
\[
 K(x,Y)\leq M\qquad\text{对所有 }x\in X,
\]
称$K$具有统一有限质量界。
逐\term{点质量}$K(x,Y)<\infty$并不保证存在共同上界，两者不能混用。
若还满足
\[
 K(x,Y)=1\qquad\text{对所有 }x,
\]
称$K$为\term{归一化测度核}（马尔可夫核）；这一附加条件用于保持总质量和边缘。

对可测映射$T:X\to Y$，取
\[
 K_T(x,B)=\mathbf 1_B(Tx)=\delta_{T(x)}(B).
\]
每个$x$对应点质量，每个$B$对应可测的示性函数，所以它是归一化测度核。
确定推前因此成为测度核传递的特例；一般核允许一项输入测度在多个目标点或集合之间分配质量，其总质量还可以随$x$变化。

\subsubsection{核积分与核复合}

\paragraph{输出测度与观测函数的传递}

给定$X$上的有限非负测度$\mu$和满足统一质量界$K(x,Y)\leq M<\infty$的核$K$，定义输出测度
\[
 (\mu K)(B)=\int_X K(x,B)\,\mu(dx).
\]
对两两不交的$B_n$，用核在固定$x$处的可数可加性和非负单调收敛，可把求和移出积分：
\[
 (\mu K)\!\left(\bigcup_nB_n\right)
  =\int_X\sum_nK(x,B_n)\,\mu(dx)
  =\sum_n\int_XK(x,B_n)\,\mu(dx),
\]
故$\mu K$是测度。
总质量满足
\[
 (\mu K)(Y)=\int_XK(x,Y)\,d\mu\leq M\mu(X).
\]
归一化核保证$(\mu K)(Y)=\mu(X)$，一般核不保证质量守恒。

对非负可测$g:Y\to[0,\infty]$，定义
\[
 Kg(x)=\int_Y g(y)\,K(x,dy).
\]
$g$为指标时，$Kg$就是核定义中的可测函数；简单函数对应有限线性组合；非负单调逼近保证$Kg$可测，并得到核积分恒等式
\begin{equation}
 \int_Y g\,d(\mu K)=\int_X Kg\,d\mu,
 \label{eq:measure-kernel-integral}
\end{equation}
允许两边共同为$+\infty$。
这里必须用到核随输入可测这第二项条件，不能只说每个截面的积分存在。
对有界可测实函数$g$，
\[
 |Kg(x)|\leq M\|g\|_\infty,
\]
各点积分有限。
一般有符号$g$在$\int|g|\,d(\mu K)<\infty$时，先对$|g|$使用式~\eqref{eq:measure-kernel-integral}，得到
\[
 \int_Y|g(y)|\,K(x,dy)<\infty
\]
在$\mu$几乎处处成立；在剩余零集上为$Kg$选定可测值，再合并正负部分即可，避免无穷相减。

用一个离散例复算。
输入点为$a,b$，$\mu(\{a\})=1$、$\mu(\{b\})=2$；输出点为$0,1$，核两行分别为$(2,1)$和$(0,2)$。
输出在$0,1$处的质量分别为
\[
 1\cdot2+2\cdot0=2,\qquad 1\cdot1+2\cdot2=5,
\]
总质量为$7$，等于按输入权重累计的行质量$1\cdot3+2\cdot2=7$，但不等于输入总质量$3$。
取$g(y)=y$，两行核积分分别为$1$和$2$，输入积分为$1\cdot1+2\cdot2=5$，与输出积分一致。

\paragraph{连续两次传递的复合}

给定核$K:X\leadsto Y$与$L:Y\leadsto Z$，统一质量界分别为$M,N<\infty$。
前一个核为每个$x$给出$Y$上的测度，后一个核以$Y$中的点为输入，按前一测度累计后一核的输出，定义
\[
 (KL)(x,C)=\int_Y L(y,C)\,K(x,dy).
\]
对固定$C$，被积函数关于$y$可测；对不交的$C_n$使用非负单调收敛得到可数可加。
质量界为
\[
 (KL)(x,Z)\leq N\,K(x,Y)\leq NM,
\]
所以复合核仍在统一有限质量界范围内；两核都归一化时复合也归一化，一般情况下不能把总质量写成$1$。

先对目标指标使用核积分恒等式，得到测度恒等式
\[
 (\mu K)L=\mu(KL);
\]
再对简单及非负函数推广，得到逐点的扩展积分恒等式
\[
 (KL)g=K(Lg).
\]
有符号$g$在$\int|g|\,d[\mu(KL)]<\infty$时，按前述方式选取可测积分版本，观测恒等式在$\mu$几乎处处意义下成立。
两个式子分别说明：先后传递输出测度与先后传回观测是一致的操作。
把$K,L$依次取为可测映射$T,S$的点核时，
\[
 K_TK_S=K_{S\circ T},
\]
恢复推前的复合式，说明按积分定义的核复合扩展了点值映射的普通复合。
本章只建立有限串联的通用运算，不展开无限串联或路径空间理论。

\subsubsection{测度值映射的可测性}

\paragraph{用集合质量检验测度族的可测依赖}

记$(Y,\mathcal B)$上全部非负测度组成的集合为$\mathcal M_+(Y)$，在其上取由评价映射
\[
 e_B(\nu)=\nu(B),\qquad B\in\mathcal B
\]
生成的$\sigma$代数，评价值取$[0,\infty]$。
在这个结构下，映射$F:X\to\mathcal M_+(Y)$可测，当且仅当对每个$B\in\mathcal B$，函数
\[
 x\longmapsto F(x)(B)
\]
都可测。
测度核的两项条件正好使$F(x)=K_x$成为这样一项测度值可测映射；统一有限质量界和归一化只进一步限制值域，不改变可测性判断。

以$Y^n$的乘积可测结构为输入，固定非负有限权重$a_1,\ldots,a_n$，定义
\[
 F(y_1,\ldots,y_n)=\sum_{i=1}^na_i\delta_{y_i}.
\]
对任意$B$，评价为$\sum_i a_i\mathbf 1_B(y_i)$，是有限个坐标可测函数之和，因此加权点测度关于输入可测，总质量有共同界$\sum_i a_i$。
取$a_i=1/n$得到\term{经验测度}这一特例。
这里要区分两个层面：每个输入给出一项测度只检查了目标端；输入到测度的映射可测还要检查这些评价函数。
当前直接指定评价$\sigma$代数，不把它无条件等同于任意空间上弱拓扑的Borel结构。
完成从点值映射到测度值映射的可测接口后，就可以形成联合测度，再判断反向分解是否存在。

\subsection{联合测度与测度\term{核分解}}

输出测度只保留目标端质量；若还要记录哪部分输入贡献了哪些输出，应保留联合测度。
本小节分别回答：给定输入测度和核时怎样构造联合权重；给定联合权重时何时能反求核。
正向构造沿用有限输入与统一质量界，反向分解则从第一边缘出发，需要额外的空间条件。
\term{条件积分}只作为理解反向分解的通用基础，其进一步代数法则另行展开。

\subsubsection{由输入测度和核构造联合测度}

\paragraph{联合权重、两个边缘与乘积特例}

给定$X$上的有限非负测度$\mu$和具有统一有限质量界的核$K:X\leadsto Y$。
先在可测矩形上规定
\[
 J(A\times B)=\int_A K(x,B)\,\mu(dx),
\]
再对乘积可测集合$E\subseteq X\times Y$用
\[
 J(E)=\int_X\left[\int_Y\mathbf 1_E(x,y)\,K(x,dy)\right]\mu(dx)
\]
构造联合测度。
由矩形指标进入，利用核的统一有限质量界把有界函数的结论经单调类推广，得到截面积分关于$x$可测；再以单调收敛核对可数可加。
总质量受$M\mu(X)$控制。
这个构造没有假定$J$在某个乘积参考测度下有密度。

在矩形上分别取$B=Y$与$A=X$，得到两个边缘：
\[
 (\pi_X)_\#J(A)=\int_A K(x,Y)\,d\mu(x),
 \qquad
 (\pi_Y)_\#J=(\mu K).
\]
第一边缘是在$\mu$上乘以行总质量$K(x,Y)$得到的测度；核归一化时它等于原输入$\mu$，一般情形不保证。
对非负乘积可测$h$，沿指标、简单函数和单调逼近有
\begin{equation}
 \int_{X\times Y}h\,dJ
  =\int_X\int_Y h(x,y)\,K(x,dy)\,\mu(dx);
 \label{eq:measure-joint-integral}
\end{equation}
有符号情形要求$\int|h|\,dJ<\infty$。
记号$\mu(dx)K(x,dy)$表示这条联合积分规则。

两个特例值得写出。
若$K(x,\cdot)=\nu$是与$x$无关的固定有限测度，矩形质量变为$\mu(A)\nu(B)$，由矩形生成族上的唯一性得$J=\mu\otimes\nu$，第一边缘为$\nu(Y)\mu$。
若$K=K_T$是可测映射$T$的归一化点核，联合测度是\term{图测度}
\[
 J=(\operatorname{id},T)_\#\mu,
\]
第一边缘为$\mu$。
乘积测度和确定图测度都是联合构造的特例，但其边缘要按各自总质量核对；反向寻找描述给定联合关系的核，需要下面的测度分解。

\subsubsection{从联合测度反求测度核}

\paragraph{子$\sigma$代数上的条件积分与观测值表示}

反向分解要先理解“保留哪些输入信息”如何表示积分。
固定有限测度空间$(\Omega,\mathcal F,\mu)$、子$\sigma$代数$\mathcal G\subseteq\mathcal F$和可积实函数$h$。
只在$\mathcal G$的集合上考察
\[
 \nu(A)=\int_A h\,d\mu,\qquad A\in\mathcal G,
\]
得到相对于\term{限制测度}$\mu|_{\mathcal G}$绝对连续的有限\term{有符号测度}；限制可测集合族保持底集不变，不能说成把$\Omega$缩小。

用$h^+$和$h^-$定义两个受限的有限非负测度
\[
 \nu_\pm(A)=\int_A h^\pm\,d\mu,
\]
它们都对$\mu|_{\mathcal G}$绝对连续。
对两者分别应用非负版本的Radon–Nikodym定理，取得$\mathcal G$可测、可积的密度$H_\pm$，令$H=H_+-H_-$的有限可测版本，则
\[
 \int_A H\,d\mu=\int_A h\,d\mu\qquad(A\in\mathcal G),
\]
并且$H$在$\mu$几乎处处意义下唯一。
这个条件积分保存每个可观测集合上的累计，不要求$H=h$逐点成立；有限测度使受限参考仍有限，不无条件推广到任意$\sigma$有限空间的子$\sigma$代数。

若$\mathcal G=\sigma(S)$由可测观测$S:\Omega\to X$生成，实值$\mathcal G$可测函数可以写成$H=\phi\circ S$，其中$\phi:X\to\mathbb R$可测。
理由是：简单的$\mathcal G$可测函数的层集合是$S$的原像，可写成$\phi(S)$；一般实值可测函数由简单函数逐级逼近，在目标端不收敛处选定一个可测值即可。
因此“依赖观测的信息”能表示成观测值的函数，而不是在每条纤维$S^{-1}(x)$上任意取值。

把这两步用于非零有限联合测度$J$：取$\Omega=X\times Y$、$\mu=J$、$S=\pi_X$、$\mathcal G=\sigma(\pi_X)$，对每个目标集合$B$令$h(x,y)=\mathbf 1_B(y)$。
条件积分因子化为$H_B=k_B\circ\pi_X$；用第一边缘$\mu_X=(\pi_X)_\#J$的输出积分公式，得到
\[
 \int_A k_B\,d\mu_X=J(A\times B)\qquad(A\in\mathcal A).
\]
右端介于$0$和$\mu_X(A)$之间，所以可选取$[0,1]$值的可测代表$k_B$，它说明输入端的条件质量函数怎样取得。
但逐个集合取得$k_B$，还没有保证这些代表在同一组输入上组成可数可加、总质量为$1$的测度；把一族条件积分协调成归一化测度核，需要下一目在标准Borel范围内的存在定理。

\paragraph{标准Borel范围内的联合分解}

反向分解当前采用两侧都是\term{标准Borel空间}的充分范围，不宣称这是最弱条件。
标准Borel结构提供可数生成族、决定测度以及协调各版本所需的条件，它没有给定跨空间距离。

\begin{theorem}[测度的核分解]
\label{thm:measure-disintegration}
设$X,Y$为标准Borel空间，$J$是$X\times Y$上的非零有限测度，第一边缘为
\[
 \mu_X=(\pi_X)_\#J.
\]
则存在归一化测度核$K:X\leadsto Y$，使
\[
 J(A\times B)=\int_A K(x,B)\,\mu_X(dx),
\]
从而
\[
 J=\mu_X(dx)\,K(x,dy)
\]
且第二边缘为$\mu_XK$。
\end{theorem}

归一化是“按实际第一边缘分解”带来的，不要求$J$的总质量为$1$；也不能反推任意给定输入和核时第一边缘必然不变。
存在性证明需要把可数决定族上的条件积分协调成同一个测度版本，本章陈述准确结论并说明这一构造障碍，完整证明另行展开。

核的唯一性是共同$\mu_X$零集以外的测度相同：若$K,L$都给出该分解，则存在同一个可测集合$N$满足$\mu_X(N)=0$，使对每个$x\notin N$和全部$B\in\mathcal B$都有
\[
 K(x,B)=L(x,B).
\]
标准Borel的可数决定族先允许合并可数个零例外，再由有限测度的唯一性推广到全部可测集合；不能把“逐集合几乎处处相等”直接写成无条件的共同零集结论。

取$X=Y=[0,1]$的Borel空间、输入测度为区间上的勒贝格测度$\lambda$、$T(x)=x$，联合图测度
\[
 J=(\operatorname{id},T)_\#\lambda
\]
集中在对角线上。
它的两个边缘都是$\lambda$，但$(\lambda\otimes\lambda)$给对角线赋零而$J$给对角线质量$1$，所以$J$相对于乘积参考没有密度；核$K(x,\cdot)=\delta_x$仍逐项重建$J$。
这个例子区分“没有联合密度”和“没有测度核分解”，也说明坐标分解与上一节相对于参考测度的\term{Lebesgue分解}回答的是不同问题。

\subsection{建立共同空间后的测度比较}

本小节先确定两项输出在哪里比较，再判断比较结果保留了哪些输入区别。
给定目标映射或共同表示，都是让同空间指标具有合法输入的方式；对输入差异的控制则以同一输出规则为前提。
联合和\term{耦合}构造用来说明可见输出与原输入的关系，不能把表示相同误写成原对象相同。

\subsubsection{目标空间与共同表示}

\paragraph{指定映射后的目标端差异}

给定$X$上的概率$P$、$Y$上的概率$Q$和预先确定的可测映射$T:X\to Y$，先形成输出测度
\[
 R=T_\#P,
\]
再在$Y$上比较$R$与$Q$。
同一个目标事件$B$对应源事件$T^{-1}(B)$，所以比较的是映射后输出对目标事件的质量判断，不能直接对原来不同空间的$P,Q$写密度相减。
若给定的是归一化测度核，则以$R=PK$形成相同的目标端问题；一般核可能改变总质量，必须先按所选比较方式核对输出范围。

若采用TV或Hellinger，在$Y$上取共同有限参考$\lambda_Y=R+Q$即可得到两项\term{Radon–Nikodym密度}，不必预设勒贝格密度。
若采用$D_{\mathrm{KL}}(R\|Q)$，要固定方向，检查$R\ll Q$和对数比的可积性；绝对连续不足以保证KL有限。
零集关系、Lebesgue分解和测试函数比较都在同一个目标可测空间上判定。

若采用Wasserstein比较，要在Polish目标空间$(Y,d_Y)$的Borel结构上取$R,Q$，固定$1\leq p<\infty$，并要求$R,Q\in\mathcal P_p(Y)$。
确定映射时，$R$的矩由输出积分写为
\[
 \int_Y d_Y(y,y_0)^p\,dR(y)
  =\int_X d_Y(Tx,y_0)^p\,dP(x);
\]
归一化测度核时则写为
\[
 \int_X\int_Y d_Y(y,y_0)^p\,K(x,dy)\,P(dx).
\]
源空间在另一个距离下的有限矩不能无条件替代目标矩。
例如以$[0,1]$上的均匀概率为输入，令
\[
 T(0)=0,\qquad T(x)=\frac1x\quad(x>0),
\]
输入在原距离下具有所有有限阶矩，但对任意$1\leq p<\infty$，
\[
 \int_0^1 T(x)^p\,dx=\int_0^1x^{-p}\,dx=+\infty,
\]
输出不属于Wasserstein有限距离的范围。

\paragraph{在共同表示空间中比较}

两侧都要重新表示时，固定可测映射$S:X\to Z$和$T:Y\to Z$，比较$S_\#P$与$T_\#Q$。
先要说明$Z$的可测结构以及两个映射各保留了什么输入区别，随后所有事件、密度和测试函数都取在这个共同空间内。
测度值规则也可由两侧各自到$Z$的核形成输出；沿用当前的归一化指标时，须核对两核都归一化；使用Wasserstein时还要检查输出矩条件，不能由一般核的可测性自动推出。

对权重差异，在$Z$上取共同参考$S_\#P+T_\#Q$，回用TV、Hellinger或指定方向的KL。
对几何差异，取$Z$为Polish度量空间$(Z,d_Z)$，要求两个表示对其Borel结构可测，固定$1\leq p<\infty$，检查两侧推前都有有限$p$阶距离矩，再计算
\[
 W_p^Z(S_\#P,T_\#Q).
\]
$X,Y$上各自的距离$d_X,d_Y$有定义，并没有给出跨空间的$d(x,y)$。
取$Z=Y$、$S$为源到目标的映射、另一侧取恒等映射，就恢复直接的目标端比较。
一般共同表示的数值依赖$Z,S,T$和所选指标；若把两侧表示都取成同一个常值，所有输入都落到同一目标点，零差异不能说明原对象相同。

\subsubsection{输出规则与可辨认范围}

\paragraph{确定映射下的信息丢失}

在同一源空间$\mathbb R$上取$P_1=\delta_{-1}$、$P_2=\delta_1$，用$T(x)=|x|$映到$[0,\infty)$。
两个源测度在事件$\{-1\}$上的质量不同，但推前后都是$\delta_1$，所以目标端的TV、Hellinger和可用的Wasserstein差异都为零；输出相同只说明这个表示看不到符号区别。

对同一可测$T$和同一源空间上的任意概率$P_1,P_2$，逐个目标事件有
\[
 |(T_\#P_1)(B)-(T_\#P_2)(B)|
  =|P_1(T^{-1}B)-P_2(T^{-1}B)|
  \leq\operatorname{TV}(P_1,P_2).
\]
对$B$取上确界得到TV在可测推前下不增：
\[
 \operatorname{TV}(T_\#P_1,T_\#P_2)\leq\operatorname{TV}(P_1,P_2),
\]
因为只检测由目标原像组成的那部分源事件。

两侧使用不同表示时不能直接套用这个估计。
即使源测度都是$\delta_0$，取$S(0)=0$、$T(0)=1$也得到$\delta_0$与$\delta_1$，TV为$1$。
同一映射的非扩张性不能改写成任意两个表示的非扩张性；粗观测也不能反向控制全部输入事件。

\paragraph{同一归一化测度核下的TV控制}

固定同一源空间上的概率$P,Q$和同一个到$Y$的归一化测度核$K$，于是$PK,QK$仍是概率。
对每个目标事件$B$，输出差为
\[
 |(PK)(B)-(QK)(B)|
  =\left|\int_XK(x,B)\,d(P-Q)(x)\right|.
\]
函数$x\mapsto K(x,B)$可测，且由核的非负性和归一化有$0\leq K(x,B)\leq1$。
回用TV关于$[0,1]$值测试函数的精确公式，逐事件得到该差不超过$\operatorname{TV}(P,Q)$，再对$B$取上确界：
\[
 \operatorname{TV}(PK,QK)\leq\operatorname{TV}(P,Q).
\]
若只调用$\|f\|_\infty\leq1$的一般有界函数界会多出系数$2$；这里正好可用核函数的非负和单位上界取得精确常数。
点核特例恢复确定映射的估计。
一般核可能改变总质量，以上归一化结论不直接适用；把一个归一化核乘常数$2$会使输出总质量加倍，已不在概率TV范围内。
两侧改用不同核时也要另行控制输出规则差异，前面的两个不同常值映射已经给出反例。
本处只证明同核TV控制，不宣布所有散度和几何指标在任意核下都收缩；KL的数据处理解释留给概率分布分析。

\paragraph{共同输入构造耦合上界}

为把输出表示连到运输比较，可取共同可测空间$(\Omega,\mathcal F)$上的概率$R$和可测映射$F:\Omega\to X$、$G:\Omega\to Y$。
$(F,G)$对乘积$\sigma$代数可测，联合推前
\[
 \gamma=(F,G)_\#R
\]
的两个边缘分别是$F_\#R$和$G_\#R$，所以它给出两项输出的一个耦合。

若两项输出都在同一个Polish度量空间$(Z,d)$中，固定$1\leq p<\infty$且两边都有有限$p$阶矩，把这个可行耦合代入Wasserstein定义：
\[
 W_p^p(F_\#R,G_\#R)
  \leq\int_\Omega d(F(\omega),G(\omega))^p\,R(d\omega).
\]
中间由推前积分计算该耦合的成本，再由下确界不超过任一候选成本得到上界。
共同输入在应用中可以是共享噪声，但这里只需可测映射和给定的$R$；该构造给出可行候选和成本上界，不保证最优，也不表明任意给定的两边缘都自带这种表示。

\subsection{给定跨空间成本的运输比较}

本小节保持源和目标各自的空间，直接规定一对源点、目标点配对要付出什么成本。
共同空间比较给出输出端差异；跨空间运输把质量守恒作为约束，在所有允许对应中累计并优化给定成本。
两者不能只因为都出现映射或耦合就等同。

\subsubsection{乘积空间上的耦合与成本}

\paragraph{跨空间质量匹配与成本累计}

把$X\times X$上的耦合推广到$X\times Y$：给定$X,Y$上的概率$P,Q$，以$\Pi(P,Q)$表示第一边缘为$P$、第二边缘为$Q$的乘积空间概率$\gamma$，即
\[
 \gamma(A\times Y)=P(A),\qquad
 \gamma(X\times B)=Q(B).
\]
两个约束分别固定出发质量与到达质量，不要求两个底层空间相同。

另行给定对乘积结构可测的非负成本函数
\[
 c:X\times Y\to[0,\infty],
\]
定义最优运输成本
\[
 \mathcal C_c(P,Q)
  =\inf_{\gamma\in\Pi(P,Q)}
    \int_{X\times Y}c(x,y)\,d\gamma(x,y).
\]
非负累计允许取$+\infty$；$c$规定的是跨空间配对代价，并不是自动存在的一条共同距离。

当两侧在同一度量空间、$c(x,y)=d(x,y)^p$，并满足前面$1\leq p<\infty$和有限$p$阶矩条件时，这个最优成本就是$W_p^p$。
一般跨空间成本的最优值不自动满足对称性、同一性或三角不等式，也不能只开$p$次方就称为\term{Wasserstein距离}。
即使两侧恰为同一空间，取$c\equiv0$也会使不同概率之间的最优成本一律为零，已不能满足距离的分离要求。

\subsubsection{运输映射与\term{运输计划}}

\paragraph{不可拆分的映射与可拆分的计划}

若要求每个源点只送到一个目标点，候选是满足
\[
 T_\#P=Q
\]
的可测映射$T:X\to Y$，其成本为
\[
 \int_X c(x,T(x))\,dP(x),
\]
这给出\term{Monge问题}。
\term{Kantorovich问题}使用全部耦合，允许把同一源点的质量拆给不同目标；两类候选遵守同样的边缘守恒要求。
标准Borel范围内的计划还可由联合分解写成
\[
 P(dx)\,K(x,dy),
\]
这里的$K$由实际第一边缘$P$分解得到，故归一化，输出约束成为$PK=Q$；点核对应确定映射。
不能把任意未归一化的核直接当作满足双边缘守恒的运输计划。

取离散源$X=\{0\}$、目标$Y=\{-1,1\}$，幂集可测，令
\[
 P=\delta_0,\qquad Q=\tfrac12\delta_{-1}+\tfrac12\delta_1.
\]
确定映射只能把全部源质量送到单个$T(0)$，没有可行运输映射；计划
\[
 \gamma=\tfrac12\delta_{(0,-1)}+\tfrac12\delta_{(0,1)}
\]
则满足两个边缘约束。
若$c(0,-1)=c(0,1)=1$，该计划成本为$1$；其核在唯一源点把质量各分一半，具体说明计划放宽了什么限制。
对任何可行映射$T$，图计划$\gamma_T=(\operatorname{id},T)_\#P$具有所需边缘，推前积分给出$\int c\,d\gamma_T=\int c(x,T(x))\,dP$，所以映射候选嵌入计划候选；并非每个计划都来自确定映射，最优计划存在也不替代最优映射存在。

\subsubsection{成本判定与比较范围}

\paragraph{可行、有限与达到分别判断}

产品耦合$P\otimes Q$始终是一个耦合，所以Kantorovich计划的可行域非空；这与原子分裂例中Monge映射可行域为空形成对照。
非空只回答是否存在守恒配对，不回答配对成本是多少。

要得到有限最优成本，至少要有一个有限成本候选。
例如
\[
 \int c\,d(P\otimes Q)<\infty
\]
足以提供这样的候选，但不是唯一途径；若$c$处处为$+\infty$，所有耦合的成本都是无穷。
有限成本条件与概率总质量为$1$是两件事，后者不能控制任意给定$c$的累计。

最小值是否达到还要另查紧性和成本正则性。
对两个\term{Polish空间}上的Borel概率，在非负\term{下半连续}成本及存在有限成本候选等标准条件下，可调用最优计划存在结论；一般可测成本不能无条件套用。
紧性怎样提供\term{弱收敛}子列、下半连续成本怎样给出积分下界，将在\ref{sec:measure-tightness}一目补齐；这里不由计划可行性推出映射存在。

\paragraph{目标匹配、运输代价与结构比较边界}

用单点源和单点目标区分两类比较。
取$X=\{a\}$、$Y=\{b\}$、$P=\delta_a$、$Q=\delta_b$、$T(a)=b$，并令$c(a,b)=7$。
目标端$T_\#P=Q$，同空间差异为零，但唯一跨空间配对的运输成本为$7$。
输出匹配问的是是否得到所需权重，运输成本问的是实现这种对应当价多大；前者为零不意味着后者很小。

因此比较前必须先给定映射、核、共同表示或跨空间成本；仅有$(X,d_X)$和$(Y,d_Y)$并没有规定跨空间点对代价$c(x,y)$。
本章按这些已经给定的对应进行测度比较，还没有从两侧内部距离自行构造空间之间的内在结构距离；这种结构比较的进一步理论另行建立。
